\documentclass[11pt]{article}
\usepackage[utf8]{inputenc}
\usepackage{geometry}
\geometry{a4paper, margin=1in}
\usepackage{enumitem}
\usepackage{hyperref}
\usepackage{amsmath}

\title{Weekly Progress Report}
\author{Pranish Uprety \\ Advisor: Dr. Sushil Poudel}
\date{Week of June 3--10, 2026}

\begin{document}
\maketitle

\section{Summary}
Last week's feedback was that the prototype is promising but the project now
needs \emph{evidence}, not more model surface. This week was a lighter week
(travel over the weekend), so I focused it on the right thing: turning
``we need evidence'' into a concrete, falsifiable experiment that the rest of
the project can be built around. The deliverable below is the evidence and
validation plan, plus the first test I propose to run.

\section{Done this week}
\begin{itemize}[leftmargin=*]
  \item Reframed the central claim as a testable hypothesis (Section 4) so it
        can be confirmed or rejected on real data, rather than asserted.
  \item Specified the exact data, unit of analysis, method, baseline, and
        success metric for the first evidence experiment.
  \item Identified a published real-world counterfactual to validate against
        (Richmond Fed 2025 difference-in-differences shutdown study), so our
        evaluation in Phase 2.4 is anchored to an external result.
  \item Confirmed the closest prior work (CausalNet, 2024) uses Granger
        causality only, which is statistical, not interventional. This keeps
        our intervention contribution distinct.
\end{itemize}

\section{Did not get to (and why)}
\begin{itemize}[leftmargin=*]
  \item No data pull or coding: I was travelling over the weekend.
  \item Carrying these to next week: download the BTS on-time slice and run
        the propagation regression in Section 4.
\end{itemize}

\section{Evidence plan: does delay actually propagate through the graph?}
The whole project rests on one premise: a disruption at one airport raises
delays at connected airports later in the day. Before modeling interventions,
I want direct evidence that this propagation is real and measurable.

\paragraph{Hypothesis (H1).} A departure-delay shock at a hub airport is
associated with elevated arrival delays at network-adjacent airports within a
lag window, beyond what local weather alone explains.

\paragraph{Data.}
\begin{itemize}[leftmargin=*]
  \item BTS Reporting Carrier On-Time Performance, 2023--2024 (delay minutes,
        origin, destination, scheduled and actual times).
  \item NOAA METAR weather per airport-hour (control variable).
  \item OpenSky for traffic density (already wired into the prototype).
\end{itemize}

\paragraph{Method.} Aggregate to the airport-hour level. For each airport,
regress its arrival delay at hour $t$ on the delays of its upstream neighbors
at hours $t-1, t-2, \dots$, controlling for local weather and time-of-day:
\[
D_{j,t} = \alpha + \sum_{k} \beta_k \, D_{\text{up}(j),\,t-k} + \gamma\, W_{j,t} + \varepsilon_{j,t}.
\]

\paragraph{Baseline.} A local-only model that uses just airport $j$'s own
history and weather, with no neighbor (graph) information.

\paragraph{Success metric.} The incremental predictive gain (for example,
$\Delta R^2$ or reduction in MAE) from adding upstream-neighbor delays. A
positive, significant $\beta_k$ that decays with graph distance and lag would
be the first concrete evidence that the graph carries delay, which is the
premise the intervention model depends on.

\paragraph{Why this matters.} If the graph adds nothing over a local model,
the GNN framing is not justified. If it does, we have measured the propagation
effect and can then ask the harder counterfactual question.

\section{Validation anchor (Phase 2.4)}
The Richmond Fed 2025 study used a difference-in-differences design on roughly
76 million flight operations to isolate about 1 million extra delay work-hours
caused by the 2025 shutdown, controlling for weather. This is an external,
peer-reviewed counterfactual estimate we can compare our model's predicted
intervention effects against, rather than grading ourselves.

\section{Questions for this meeting}
\begin{itemize}[leftmargin=*]
  \item Is H1 the right first evidence target, or would you rather I start
        from the intervention side directly?
  \item For the propagation test, do you prefer a simple lagged regression
        first, or jumping to a Granger-causality formulation like CausalNet?
  \item Which slice of BTS data should I start with: full 2023--2024, or one
        representative disruption month?
\end{itemize}

\section{Plan for next week}
\begin{itemize}[leftmargin=*]
  \item Download the BTS on-time slice and build the airport-hour table.
  \item Run the local-only baseline and the neighbor-augmented model.
  \item Report the incremental gain as the first evidence figure.
\end{itemize}

\end{document}
